\documentclass[../../main.tex]{subfiles}
\graphicspath{{\subfix{../../image/}}}

\begin{document}

\section{习题}\label{section:ch07-5-习题}

以下习题中均设 $U\subset\mathbb{R}^n$ 为开、有界集，边界光滑，$T>0$。

\begin{exercise}\label{exercise:ch07-5-1}
证明：带 Neumann 边界条件的热方程初边值问题
\begin{align*}
\begin{cases}
u_t-\Delta u=f & \text{于 } U_T,\\
\dfrac{\partial u}{\partial\nu}=0 & \text{于 } \partial U\times[0,T],\\
u=g & \text{于 } U\times\{t=0\}
\end{cases}
\end{align*}
至多存在一个光滑解。
\end{exercise}

\begin{exercise}\label{exercise:ch07-5-2}
设 $u$ 是
\begin{align*}
\begin{cases}
u_t-\Delta u=0 & \text{于 } U\times(0,\infty),\\
u=0 & \text{于 } \partial U\times[0,\infty),\\
u=g & \text{于 } U\times\{t=0\}
\end{cases}
\end{align*}
的光滑解。证明指数衰减估计：
\begin{align*}
\|u(\cdot,t)\|_{L^2(U)}\leqslant e^{-\lambda_1 t}\|g\|_{L^2(U)}\qquad (t\geqslant 0),
\end{align*}
其中 $\lambda_1>0$ 是 $-\Delta$（零边界条件）在 $U$ 上的主特征值。
\end{exercise}

\begin{exercise}[对偶动力学。]\label{exercise:ch07-5-3}设 $u$ 是
\begin{align*}
\begin{cases}
u_t+Lu=0 & \text{于 } U_T,\\
u=0 & \text{于 } \partial U\times[0,T],\\
u=g & \text{于 } U\times\{t=0\}
\end{cases}
\end{align*}
的光滑解，其中 $L$ 表示二阶椭圆算子；又设 $v$ 是对偶问题
\begin{align*}
\begin{cases}
v_t-L^*v=0 & \text{于 } U_T,\\
v=0 & \text{于 } \partial U\times[0,T],\\
v=h & \text{于 } U\times\{t=T\}
\end{cases}
\end{align*}
的光滑解。证明
\begin{align*}
\int_U v(\cdot,0)g\,\mathrm{d}x=\int_U u(\cdot,T)h\,\mathrm{d}x.
\end{align*}
\end{exercise}

\begin{exercise}[Poisson 方程的Galerkin方法。]\label{exercise:ch07-5-4}设 $f\in L^2(U)$，并设 $u_m=\sum_{k=1}^m d_m^k w_k$ 满足
\begin{align*}
\int_U Du_m\cdot Dw_k\,\mathrm{d}x=\int_U f\cdot w_k\,\mathrm{d}x
\end{align*}
对 $k=1,\cdots,m$，其中 $\{w_k\}_{k=1}^\infty$ 是 $H^1_0(U)$ 的一组正交基。证明 $\{u_m\}_{m=1}^\infty$ 在 $H^1_0(U)$ 中弱收敛于
\begin{align*}
\begin{cases}
-\Delta u=f & \text{于 } U,\\
u=0 & \text{于 } \partial U
\end{cases}
\end{align*}
的弱解 $u$。
\end{exercise}

\begin{exercise}\label{exercise:ch07-5-5}
设
\begin{align*}
\begin{cases}
u_k\rightharpoonup u & \text{弱收敛于 } L^2(0,T;H^1_0(U)),\\
u_k'\rightharpoonup v & \text{弱收敛于 } L^2(0,T;H^{-1}(U)).
\end{cases}
\end{align*}
证明 $v=u'$。（提示：取 $\phi\in C_c^1(0,T)$、$w\in H^1_0(U)$。则
\begin{align*}
\int_0^T\langle u_k',\phi w\rangle\,\mathrm{d}t=-\int_0^T\langle u_k,\phi' w\rangle\,\mathrm{d}t.)
\end{align*}
\end{exercise}

\begin{exercise}\label{exercise:ch07-5-6}
设 $H$ 是 Hilbert 空间，$u_k\rightharpoonup u$ 于 $L^2(0,T;H)$。又设有一致界
\begin{align*}
\operatorname*{ess\,sup}_{0\leqslant t\leqslant T}\|u_k(t)\|\leqslant C\qquad (k=1,\cdots)
\end{align*}
（某常数 $C$）。证明
\begin{align*}
\operatorname*{ess\,sup}_{0\leqslant t\leqslant T}\|u(t)\|\leqslant C.
\end{align*}
（提示：对 $0\leqslant a\leqslant b\leqslant T$ 及 $v\in H$，有 $\int_a^b(v,u_k(t))\,\mathrm{d}t\leqslant C\|v\||b-a|$。）
\end{exercise}

\begin{exercise}\label{exercise:ch07-5-7}
设 $u$ 是
\begin{align*}
\begin{cases}
u_t-\Delta u+cu=0 & \text{于 } U\times(0,\infty),\\
u=0 & \text{于 } \partial U\times[0,\infty),\\
u=g & \text{于 } U\times\{t=0\}
\end{cases}
\end{align*}
的光滑解，且函数 $c$ 满足 $c\geqslant\gamma>0$。证明逐点指数衰减估计
\begin{align*}
|u(x,t)|\leqslant Ce^{-\gamma t}\qquad ((x,t)\in U\times(0,\infty)).
\end{align*}
\end{exercise}

\begin{exercise}\label{exercise:ch07-5-8}
设 $u$ 是习题\cref{exercise:ch07-5-7} 中 PDE 的光滑解，$g\geqslant 0$，且 $c$ 有界（但不一定非负）。证明 $u\geqslant 0$。（提示：$v\triangleq e^{-\lambda t}u$ 满足什么 PDE？）
\end{exercise}

\begin{exercise}\label{exercise:ch07-5-9}
证明\cref{subsection:ch07-1-3-正则性} 中不等式 (54)、\cref{subsection:ch07-2-3-正则性} 中不等式\eqref{align:ch07-2-eq59}。
（提示：设 $u$ 光滑，$u=0$ 于 $\partial U$。两次分部积分以变换项 $(Lu,-\Delta u)$，再估计边界项。局部作变量替换并使用截断函数后，可设边界平坦。本题较难。）
\end{exercise}

\begin{exercise}\label{exercise:ch07-5-10}
证明电报方程初边值问题
\begin{align*}
\begin{cases}
u_{tt}+du_t-u_{xx}=f & \text{于 } (0,1)\times(0,T),\\
u=0 & \text{于 } (\{0\}\times[0,T])\cup(\{1\}\times[0,T]),\\
u=g,\ u_t=h & \text{于 } (0,1)\times\{t=0\}
\end{cases}
\end{align*}
至多存在一个光滑解。其中 $d$ 为常数。
\end{exercise}

\begin{exercise}\label{exercise:ch07-5-11}
证明梁方程问题
\begin{align*}
\begin{cases}
u_{tt}+u_{xxxx}=0 & \text{于 } (0,1)\times(0,T),\\
u=u_x=0 & \text{于 } (\{0\}\times[0,T])\cup(\{1\}\times[0,T]),\\
u=g,\ u_t=h & \text{于 } [0,1]\times\{t=0\}
\end{cases}
\end{align*}
至多存在一个光滑解 $u$。
\end{exercise}

\begin{exercise}\label{exercise:ch07-5-12}
证明\cref{subsection:ch07-4-1-定义与初等性质} 中的预解恒等式\eqref{align:ch07-4-eq12} 与\eqref{align:ch07-4-eq13}。
\end{exercise}

\begin{exercise}\label{exercise:ch07-5-13}
当 $u\in D(A)$ 时，证明\cref{subsection:ch07-4-1-定义与初等性质} 公式\eqref{align:ch07-4-eq16} 中所使用的等式
\begin{align*}
A\int_0^\infty e^{-\lambda t}S(t)u\,\mathrm{d}t=\int_0^\infty e^{-\lambda t}AS(t)u\,\mathrm{d}t
\end{align*}
成立。（提示：用 Riemann 和逼近积分，并回忆 $A$ 是闭算子。）
\end{exercise}

\begin{exercise}\label{exercise:ch07-5-14}
对 $t>0$ 定义
\begin{align*}
[S(t)g](x)=\int_{\mathbb{R}^n}\Phi(x-y,t)g(y)\,\mathrm{d}y\qquad (x\in\mathbb{R}^n),
\end{align*}
其中 $g:\mathbb{R}^n\to\mathbb{R}$，$\Phi$ 是热方程的基本解。又令 $S(0)g=g$。

(a) 证明 $\{S(t)\}_{t\geqslant 0}$ 是 $L^2(\mathbb{R}^n)$ 上的压缩半群。

(b) 证明 $\{S(t)\}_{t\geqslant 0}$ 不是 $L^\infty(\mathbb{R}^n)$ 上的压缩半群。
\end{exercise}

\begin{exercise}\label{exercise:ch07-5-15}
设 $\{S(t)\}_{t\geqslant 0}$ 是 $X$ 上的压缩半群，生成元为 $A$。归纳地定义
\begin{align*}
D(A^k)\triangleq \{u\in D(A^{k-1})\mid A^{k-1}u\in D(A)\}\qquad (k=2,\cdots).
\end{align*}
证明：若对某个 $k$ 有 $u\in D(A^k)$，则对每个 $t\geqslant 0$，$S(t)u\in D(A^k)$。
\end{exercise}

\begin{exercise}\label{exercise:ch07-5-16}
利用习题\cref{exercise:ch07-5-15} 证明：若 $u$ 是 $X=L^2(U)$ 中
\begin{align*}
\begin{cases}
u_t-\Delta u=0 & \text{于 } U_T,\\
u=0 & \text{于 } \partial U\times[0,T],\\
u=g & \text{于 } U\times\{t=0\}
\end{cases}
\end{align*}
的半群解，且 $g\in C_c^\infty(U)$，则对每个 $0\leqslant t\leqslant T$，$u(\cdot,t)\in C^\infty(U)$。
\end{exercise}

\end{document}
